% !TEX root = ../notes_template.tex
\chapter{Elbow, Forearm, Wrist \& Hand}\label{chp:distal_ue}

\section*{Objectives}
\begin{itemize}
  \item Perform and interpret MMT for elbow, forearm, wrist, hand, and thumb motions.
  \item Identify innervation and nerve root levels for each tested muscle.
  \item Conduct MLT for elbow and wrist flexor/extensor groups, including biceps and triceps brachii.
  \item Grade and interpret DTRs for the biceps and triceps tendons.
  \item Palpate and image the brachial artery using ultrasound.
  \item Observe and trace common upper extremity nerve entrapment sites using anatomical landmarks.
\end{itemize}

\section*{Peripheral Nerve Entrapment and Screening Strategy}

\noindent
Peripheral nerves can become compressed or irritated along their course, producing motor and sensory impairments distal to the entrapment. In this lab, we focus on screening for **median, ulnar, and radial** nerve integrity by combining anatomical palpation with MMT and reflex testing.

\vspace{0.5cm}

We adopt a **proximal-to-distal logic** when tracing nerves, as this mirrors both clinical reasoning and the path of nerve vulnerability. Entrapment sites are reviewed below with location cues. You will relate these sites to your testing results throughout the elbow, forearm, and hand.

\vspace{0.5cm}

\begin{table}[H]
\centering
\begin{tabular}{|p{2.8cm}|p{7cm}|p{4.7cm}|}
\hline
\textbf{Nerve} & \textbf{Entrapment Site} & \textbf{Memory Cue} \\
\hline
\multirow{4}{*}{\textbf{Median}} 
  & Bicipital Aponeurosis (anterior elbow) & “Elbow pit tunnel” \\
  & Lig of Struthers (medial distal humerus) & “Hidden elbow bridge” \\
  & Pronator Teres (proximal forearm) & “Forearm muscle trap” \\
  & Carpal Tunnel (wrist, flexor retinaculum) & “Classic compression” \\
\hline
\multirow{4}{*}{\textbf{Radial}} 
  & Radial Groove (posterior humerus) & “Humeral hug” \\
  & Radial Head (lateral elbow) & “Turning the radius” \\
  & Arcade of Frohse (edge of supinator) & “Tight supinator turn” \\
  & Supinator (deep in forearm) & “Radial squeeze” \\
\hline
\multirow{3}{*}{\textbf{Ulnar}} 
  & Arcade of Struthers (medial arm) & “Pre-cubital trap” \\
  & Cubital Tunnel (post medial epicondyle) & “Funny bone spot” \\
  & Flexor Carpi Ulnaris (between heads) & “Final tunnel” \\
\hline
\end{tabular}
\captionsetup{justification=raggedright, singlelinecheck=false}
\caption{Peripheral Nerve Entrapment Sites of the Upper Extremity (Proximal to Distal). \textit{Note: Identifying the specific location of nerve entrapment is beyond the scope of this lab and will not be assessed on the practical exam. This table is included to introduce the anatomical reasoning that underlies neuromuscular testing and to support clinical pattern recognition over time.}}
\end{table}

\section*{Elbow and Forearm Movements: Muscles and Innervation}

\vspace{.5cm}

\begin{tabular}{|p{4cm}|p{4.5cm}|p{4cm}|p{2.5cm}|}
\hline
\textbf{Movement} & \textbf{Muscle} & \textbf{Peripheral Nerve} & \textbf{Nerve Root(s)} \\
\hline
\multirow{3}{*}{Elbow Flexion}
  & Biceps Brachii & Musculocutaneous & C5–C6 \\
  & Brachialis & Musculocutaneous & C5–C6 \\
  & Brachioradialis & Radial & C5–C6 \\
\hline
Elbow Extension
  & Triceps Brachii & Radial & C6–C8 \\
\hline
Forearm Supination
  & Biceps Brachii & Musculocutaneous & C5–C6 \\
  & Supinator & Radial (deep branch) & C6–C7 \\
\hline
Forearm Pronation
  & Pronator Teres & Median & C6–C7 \\
  & Pronator Quadratus & Median  & C7–C8 \\
\hline
\end{tabular}

\begin{labactivitybox}
\textbf{Activity: MMT – Elbow and Forearm Movements}

\begin{itemize}
  \item Perform gravity-resisted and gravity-eliminated MMT for:
  \begin{itemize}
    \item Elbow flexors (biceps brachii, brachialis, brachioradialis)
    \item Elbow extensors (triceps brachii)
    \item Forearm supinators (biceps brachii, supinator)
    \item Forearm pronators (pronator teres, pronator quadratus)
  \end{itemize}
  \item Conduct a break test using the HHD for either elbow flexion or extension.
  \item Record stabilization strategies, resistance application, and observed compensatory motions.
  \item For each tested muscle, document its peripheral innervation and most common nerve root(s).
\end{itemize}
\end{labactivitybox}

\noindent Deep tendon reflex testing helps corroborate segmental integrity assessed through MMT and supports differential diagnosis of upper limb weakness.

\begin{labactivitybox}
\textbf{Activity: DTR Testing – Biceps and Triceps}

\begin{itemize}
  \item Perform reflex testing for:
  \begin{itemize}
    \item \textbf{Biceps tendon} – C5–C6 (Musculocutaneous nerve)
    \item \textbf{Triceps tendon} – C7 (Radial nerve)
  \end{itemize}
  \item Grade each reflex using the standard 0–4+ scale.
  \item Compare left vs. right for symmetry and amplitude.
  \item Correlate reflex findings with muscle strength results from MMT.
\end{itemize}
\end{labactivitybox}

\begin{labactivitybox}
\textbf{Activity: Muscle Length Testing – Upper Limb (Elbow)}
\begin{itemize}
  \item Muscles: Biceps brachii, Triceps brachii
  \item Document R1, R2, and classify muscle length.
\end{itemize}
\end{labactivitybox}


\section*{Wrist Movements: Muscles and Innervation}

\vspace{0.5cm}

\begin{tabular}{|p{4cm}|p{4.5cm}|p{4cm}|p{2.5cm}|}
\hline
\textbf{Movement} & \textbf{Muscle} & \textbf{Peripheral Nerve} & \textbf{Nerve Root(s)} \\
\hline
\multirow{2}{*}{Wrist Flexion}
  & Flexor Carpi Radialis & Median & C6–C7 \\
  & Flexor Carpi Ulnaris & Ulnar & C7–T1 \\
\hline
\multirow{2}{*}{Wrist Extension}
  & Ext Carpi Radialis Longus & Radial & C6–C7 \\
  & Extensor Carpi Ulnaris & Radial & C7–C8 \\
\hline
\multirow{2}{*}{Radial Deviation}
  & Flexor Carpi Radialis & Median & C6–C7 \\
  & Ext Carpi Radialis Longus & Radial & C6–C7 \\
\hline
\multirow{2}{*}{Ulnar Deviation}
  & Flexor Carpi Ulnaris & Ulnar & C7–T1 \\
  & Extensor Carpi Ulnaris & Radial & C7–C8 \\
\hline
\end{tabular}

\begin{labactivitybox}
\textbf{Activity: MMT – Forearm Rotation and Wrist Motion}

\begin{itemize}
  \item Test pronator teres, pronator quadratus, and supinator for forearm rotation in gravity-resisted and -eliminated positions.
  \item Test wrist flexors and extensors for strength, symmetry, and compensation.
  \item Repeat one test using HHD.
  \item Identify and record innervation and spinal root levels for each muscle.
\end{itemize}
\end{labactivitybox}

\section*{Hand Extrinsic Muscles (Gross Motor Screening Focus)}

\noindent
For clinical screening and gross motor assessment, we focus on muscle groups with functional or neurodiagnostic relevance. This includes composite grip function and selected muscles innervated by each major nerve (median, ulnar, radial).

\begin{table}[H]
\centering
\captionsetup{justification=raggedright, singlelinecheck=false}
\begin{tabular}{|p{4cm}|p{4.5cm}|p{4cm}|p{2.5cm}|}
\hline
\textbf{Movement} & \textbf{Muscle} & \textbf{Peripheral Nerve} & \textbf{Nerve Root(s)} \\
\hline
\multirow{3}{*}{Grip} 
  & FDS & Median & C7–T1 \\
  & FDP (lat. 2–3) & Median & C8–T1 \\
  & FDP (med. 4–5) & Ulnar & C8–T1 \\
\hline
MCP Ext (digits 2–5)
  & Extensor Digitorum & Radial & C7–C8 \\
\hline
\end{tabular}
\caption{\textit{Note: Grip strength primarily reflects extrinsic flexors but is also influenced by intrinsic hand muscles such as interossei and thenar muscles.}}
\end{table}

\vspace{0.5cm}

\begin{labactivitybox}
\textbf{Activity: MMT – Gross Motor Screening of Hand Function with Functional Grip Testing}

\begin{itemize}
  \item Perform \textbf{hand grip testing} using a dynamometer in three positions:
  \begin{enumerate}
    \item Neutral wrist
    \item Wrist extension (20–30°)
    \item Wrist flexion (45°)
  \end{enumerate}
  \item Compare force output and interpret differences using principles of:
  \begin{itemize}
    \item \textbf{Active Insufficiency} – wrist flexors (grip force ↓ in wrist flexion)
    \item \textbf{Passive Insufficiency} – wrist extensors (tension limits grip in wrist extension)
  \end{itemize}
  \item Manually test:
    \begin{itemize}
      \item \textbf{Extensor Digitorum} – radial nerve screen
      \item \textbf{FDP digits 4–5} – ulnar nerve screen
    \end{itemize}
  \item Record findings and note any weakness, asymmetry, or compensatory patterns.
  \item Document the peripheral nerve and spinal root(s) for each muscle tested.
\end{itemize}
\end{labactivitybox}


\section*{Hand Intrinsic Muscles (Screening Focus)}

\noindent
These intrinsic muscles are emphasized for clinical screening due to their role in functional grip and their diagnostic utility in assessing median and ulnar nerve integrity.

\begin{table}[H]
\centering
\captionsetup{justification=raggedright, singlelinecheck=false}
\begin{tabular}{|p{4cm}|p{4.5cm}|p{4cm}|p{2.5cm}|}
\hline
\textbf{Movement} & \textbf{Muscle} & \textbf{Peripheral Nerve} & \textbf{Nerve Root(s)} \\
\hline
Thumb Abduction
  & Abductor Pollicis Brevis & Median & C8–T1 \\
\hline
Thumb Adduction
  & Adductor Pollicis & Ulnar & C8–T1 \\
\hline
Index Finger Abd
  & 1st Dorsal Interosseous & Ulnar & C8–T1 \\
\hline
Digits 3–5 Abduction
  & Dorsal Interossei (2–4) & Ulnar & C8–T1 \\
\hline
Digits 3–5 Adduction
  & Palmar Interossei (2, 4, 5) & Ulnar & C8–T1 \\
\hline
MCP Flexion + IP Extension (Digits 2–3)
  & Lumbricals 1–2 & Median & C8–T1 \\
\hline
MCP Flexion + IP Extension (Digits 4–5)
  & Lumbricals 3–4 & Ulnar & C8–T1 \\
\hline
\end{tabular}
\caption{Movements of the Hand for Functional and Neurological Screening Involving Intrinsic Muscles}
\end{table}

\vspace{0.5cm}

\begin{labactivitybox}
\textbf{Activity: MMT – Intrinsic Muscle Screening of the Hand}

\begin{itemize}
  \item Manually test:
  \begin{itemize}
    \item \textbf{Abductor Pollicis Brevis} – median nerve screen
    \item \textbf{Adductor Pollicis} – ulnar nerve screen
    \item \textbf{1st Dorsal Interosseous} – ulnar nerve screen
    \item \textbf{Digits 3–5 Abduction/Adduction} – interossei for ulnar integrity
    \item \textbf{Lumbricals} (1–2 vs. 3–4) – differential median vs. ulnar screen
  \end{itemize}
  \item Observe for:
  \begin{itemize}
    \item Atrophy of thenar and first dorsal interosseous spaces
    \item Difficulty with fine grasp or pinch
    \item Imbalance between radial and ulnar side strength
  \end{itemize}
  \item Document peripheral nerve and nerve root(s) for each muscle tested.
\end{itemize}
\end{labactivitybox}


\section*{Neurovascular Ultrasound and Palpation of the Distal Upper Extremity}

\begin{labactivitybox}
\textbf{Activity: Vascular Assessment – Brachial and Radial Arteries}

\textit{Palpation + Ultrasound Visualization}
\begin{itemize}
  \item Locate and palpate the \textbf{brachial artery} in the antecubital fossa and the \textbf{radial artery} at the wrist.
  \item Use ultrasound in SAX to confirm vessel identity through \textbf{compressibility} and \textbf{pulsatility}.
  \item Observe for surrounding structures and note any asymmetry.
  \item Record findings, image depth, and probe orientation.
\end{itemize}
\end{labactivitybox}

\begin{labactivitybox}
\textbf{Activity: Peripheral Nerve Scanning – Median and Ulnar Nerves}

\textit{RUSI Localization and Entrapment Screening}

\begin{itemize}
  \item \textbf{Median Nerve:}
  \begin{itemize}
    \item Use \textbf{SAX} at the level of the carpal tunnel (between the scaphoid and pisiform) to identify the nerve’s “honeycomb” fascicular pattern deep to the flexor retinaculum.
    \item Confirm landmarks: radius (Rad), lunate (Lun), palmaris longus (PL), flexor tendons.
    \item Use \textbf{LAX} to trace the ribbon-like hypoechoic nerve as it courses deep to PL and over the lunate.
    \item Apply dynamic imaging (thumb opposition) to confirm nerve vs. tendon identity.
  \end{itemize}

  \item \textbf{Ulnar Nerve:}
  \begin{itemize}
    \item Use \textbf{SAX} at the cubital tunnel (between the medial epicondyle and olecranon) and at Guyon’s canal (level of pisiform).
    \item Identify key landmarks: pisiform (Pis), flexor digiti minimi brevis (FDMB), ulnar artery (UA), and flexor retinaculum (FR).
    \item Visualize the ulnar nerve as a hypoechoic bundle superficial to the retinaculum and lateral to the pisiform.
    \item Confirm normal fascicular pattern and assess for swelling or distortion.
  \end{itemize}

  \item \textbf{Interpretation:}
  \begin{itemize}
    \item Note probe orientation (SAX/LAX), anatomical level, image depth, and signs of entrapment (e.g., enlargement, compression, anisotropy).
  \end{itemize}
\end{itemize}
\end{labactivitybox}

\section*{Clinical Integration}

\begin{itemize}
  \item How did gross motor screening using grip, extrinsic, and intrinsic hand muscle tests guide your clinical reasoning about nerve involvement?
  \item In what ways did your MMT findings align with or differ from your reflex testing (biceps and triceps)?
  \item What patterns of weakness or asymmetry might suggest proximal vs. distal nerve entrapment?
  \item How did ultrasound imaging of the median and ulnar nerves enhance your understanding of nerve anatomy and potential entrapment zones?
  \item How might reflex integrity and motor strength together inform a diagnosis of cervical radiculopathy versus peripheral nerve compression?
\end{itemize}

\section*{Next Steps}

\begin{itemize}
  \item Consolidate your understanding of distal upper extremity innervation by reviewing dermatomes, myotomes, and reflexes.
  \item Practice the MMT procedures introduced this week for extrinsic grip, intrinsic hand muscles, and selected nerve screens.
  \item Practice the RUSI procedures for imaging the brachial and radial arteries, and the median and ulnar nerves in SAX and LAX views.
  \item Review the anatomy of the pelvic girdle and hip to prepare for Week 5, focusing on similarities in regional interdependence logic.
\end{itemize}

\section*{Checklist Summary}

\begin{itemize}
  \item [$\square$] Performed gravity-resisted and gravity-eliminated MMT for elbow flexion, extension, supination, and pronation with correct stabilization and substitution awareness.
  \item [$\square$] Used hand-held dynamometry (HHD) to quantify force output for either elbow flexion or extension.
  \item [$\square$] Performed and graded deep tendon reflexes for the biceps (C5–C6) and triceps (C7) tendons.
  \item [$\square$] Conducted MMT for extrinsic hand function (grip, finger extension, FDP) to screen radial and ulnar nerve involvement.
  \item [$\square$] Compared grip strength with wrist in flexed vs. extended positions to evaluate influence of active and passive insufficiency.
  \item [$\square$] Conducted MMT for intrinsic hand muscles (APB, adductor pollicis, 1st dorsal interosseous, lumbricals) for differential median vs. ulnar nerve screening.
  \item [$\square$] Identified and interpreted upper extremity nerve entrapment sites from elbow to wrist using anatomical landmarks.
  \item [$\square$] Acquired and interpreted ultrasound images of the median and ulnar nerves in both SAX and LAX views using dynamic validation techniques.
  \item [$\square$] Palpated and imaged the brachial and radial arteries using ultrasound to confirm pulsatility and compressibility.
\end{itemize}

